\section{Conclusions}

Two related issues in biodiversity monitoring were the focus of this project:
the difficulty of collecting large-scale bird observation data using conventional
methods, and the lack of accessible tools that allow non-expert volunteers to
contribute reliable species records \cite{fraisl2022citizen,johnston2022outstanding}.
To address these issues, BirdSight was developed as a citizen science platform
integrated with a machine learning model.

The project's machine learning component used a dataset of 150,000 bird photos from 30 species
that are commonly observed in Turkey to train and compare five CNN architectures. To guarantee a fair 
comparison, all models were trained under the same circumstances. The outcomes demonstrated a 
distinct performance difference across the architectures. Given its comparatively high parameter 
count and absence of residual or dense connections, VGG16 yielded the worst results. Considering 
its lightweight nature, MobileNetV2 did rather well, but its classification accuracy was not as 
good as that of the deeper models. Both ResNet50 and DenseNet121 produced results that were 
competitive, with accuracy scores of 88\%. InceptionV3 achieved 89.1\% accuracy on the original 
23-class dataset, outperforming all other architectures in every assessed metric.

After the first comparison, a number of focused experiments were conducted in an effort to enhance 
the InceptionV3 model. To improve generalisation and lessen overconfident predictions, label 
smoothing was used. To make the model more resilient to the inherent diversity found in field photos
taken by amateur birdwatchers, improved data augmentation techniques were used, such as vertical 
flipping, random erasing, and broader colour jitter ranges \cite{johnston2022outstanding}. An optimal 
freeze epoch value of 11 was identified through experimentation. When all three changes were combined, the accuracy on 
the 23-class dataset was 89.5\%. The final model maintained a consistent accuracy of 89.3\% with a 
Top-3 accuracy of 95.9\% and a Top-5 accuracy of 97.8\% when the dataset was increased to 30 species 
and 150,000 photos. These results confirm that the model can reliably identify bird species from 
amateur photographs, which is the primary requirement for integration into a citizen science platform.

Alongside the machine learning component, the BirdSight web platform was also developed.
Spring Boot was employed in the back-end of the platform’s three-layer architecture
to handle API operations and business logic. PostgreSQL was used as the database
management system, while the PostGIS extension handled geospatial observation data. 
The front-end was developed using Next.js and handled the user interface, API requests, 
and client-side functionality. MinIO is the object storage solution, and the InceptionV3 model is implemented as a REST microservice in the backend. It was developed separately from the application codebase, making it easy to upgrade and maintain. This separation of duties made it easier to update and maintain the system since it separated the model service from the application logic. 
The core functionalities of the platform such as submitting observations, receiving species
predictions, browsing the observation feed, and exploring sightings on an interactive
map were all successfully implemented.

Overall, BirdSight demonstrates that combining transfer learning with a well-designed
citizen science platform is a practical and effective approach to supporting large-scale
bird population monitoring \cite{lolge2025evolution}. The system lowers the barrier for public participation
by automating the species identification step, which is typically the main obstacle
for non-expert volunteers. At the same time, the geospatial data collected through
the platform has the potential to support ecological research by providing structured,
location-tagged observation records at a scale that traditional fieldwork methods
cannot match.

\clearpage